\documentclass[11pt]{article}
\usepackage[margin=1in]{geometry}
\usepackage{booktabs}
\usepackage{graphicx}
\usepackage{makecell}
\usepackage{caption}
\usepackage{microtype}
\usepackage[hidelinks]{hyperref}

\title{Original Assignment ATEs and Campaign-Day Take-Up}
\author{}
\date{}

\begin{document}
\maketitle

\section{Original assignment effects}

The endpoint analysis uses the originally randomized Close/Far assignment rather
than the subsequently corrected distance classification. The existing estimates
are reproduced below. Treatment rows are effects relative to Control; the Control
row reports predicted levels.

\begin{table}[htbp]
\centering
\caption{Original distance-assignment intention-to-treat estimates}
\input{../../presentations/rf-tables/main-specs/rf_original_distance_itt_tbl.tex}
\end{table}

\clearpage
\section{Campaign-day accumulation}

For campaign day $t$, cumulative take-up is an indicator for directly recorded
treatment on or before day $t$. Daily incidence is an indicator for directly
recorded treatment exactly on day $t$. Both outcomes retain the full randomized
sample; in particular, daily incidence is not conditioned on a post-treatment
risk set. Day 12 therefore equals the endpoint used in the original-assignment
ATE table exactly.

We estimate the same linear specification used for that endpoint, including the
four treatment arms interacted with original Close/Far assignment, gender, age,
expected distance, and county fixed effects. Ribbons and table intervals are
pointwise 95\% intervals from paired exponential cluster-weight draws at the 144
randomization clusters. The same weights are used across all campaign days.

\begin{figure}[htbp]
\centering
\includegraphics[width=\linewidth]{figures/campaign-day-cumulative.pdf}
\caption{Adjusted cumulative take-up over the campaign. Shaded regions are
pointwise 95\% cluster-weight intervals.}
\end{figure}

\begin{table}[htbp]
\centering
\caption{Adjusted cumulative take-up on selected campaign days}
\resizebox{\linewidth}{!}{\begin{tabular}{lcccc}
\toprule
 & Day 1 & Day 3 & Day 6 & Day 12 \\
\midrule
\multicolumn{5}{l}{\textit{Panel A: Close assignment}} \\
Control & \makecell[c]{0.111\\{[0.083, 0.139]}} & \makecell[c]{0.207\\{[0.167, 0.244]}} & \makecell[c]{0.327\\{[0.280, 0.371]}} & \makecell[c]{0.440\\{[0.389, 0.487]}} \\
Calendar & \makecell[c]{0.099\\{[0.072, 0.129]}} & \makecell[c]{0.238\\{[0.194, 0.284]}} & \makecell[c]{0.338\\{[0.289, 0.388]}} & \makecell[c]{0.456\\{[0.423, 0.493]}} \\
Ink & \makecell[c]{0.067\\{[0.044, 0.096]}} & \makecell[c]{0.145\\{[0.108, 0.187]}} & \makecell[c]{0.200\\{[0.151, 0.250]}} & \makecell[c]{0.313\\{[0.259, 0.370]}} \\
Bracelet & \makecell[c]{0.102\\{[0.073, 0.134]}} & \makecell[c]{0.242\\{[0.201, 0.291]}} & \makecell[c]{0.361\\{[0.306, 0.416]}} & \makecell[c]{0.474\\{[0.406, 0.535]}} \\
\addlinespace
\multicolumn{5}{l}{\textit{Panel B: Far assignment}} \\
Control & \makecell[c]{0.035\\{[0.024, 0.047]}} & \makecell[c]{0.090\\{[0.063, 0.119]}} & \makecell[c]{0.153\\{[0.115, 0.190]}} & \makecell[c]{0.259\\{[0.204, 0.323]}} \\
Calendar & \makecell[c]{0.027\\{[0.018, 0.036]}} & \makecell[c]{0.085\\{[0.062, 0.107]}} & \makecell[c]{0.166\\{[0.142, 0.190]}} & \makecell[c]{0.265\\{[0.232, 0.297]}} \\
Ink & \makecell[c]{0.038\\{[0.025, 0.051]}} & \makecell[c]{0.099\\{[0.068, 0.134]}} & \makecell[c]{0.169\\{[0.116, 0.226]}} & \makecell[c]{0.305\\{[0.240, 0.362]}} \\
Bracelet & \makecell[c]{0.041\\{[0.028, 0.055]}} & \makecell[c]{0.116\\{[0.091, 0.145]}} & \makecell[c]{0.210\\{[0.176, 0.242]}} & \makecell[c]{0.358\\{[0.314, 0.408]}} \\
\addlinespace
\multicolumn{5}{l}{\textit{Panel C: Far minus Close}} \\
Control & \makecell[c]{-0.076\\{[-0.105, -0.049]}} & \makecell[c]{-0.117\\{[-0.163, -0.071]}} & \makecell[c]{-0.175\\{[-0.230, -0.121]}} & \makecell[c]{-0.181\\{[-0.254, -0.098]}} \\
Calendar & \makecell[c]{-0.073\\{[-0.104, -0.043]}} & \makecell[c]{-0.154\\{[-0.204, -0.102]}} & \makecell[c]{-0.172\\{[-0.227, -0.116]}} & \makecell[c]{-0.191\\{[-0.241, -0.142]}} \\
Ink & \makecell[c]{-0.029\\{[-0.060, -0.002]}} & \makecell[c]{-0.046\\{[-0.102, 0.006]}} & \makecell[c]{-0.031\\{[-0.106, 0.044]}} & \makecell[c]{-0.008\\{[-0.091, 0.076]}} \\
Bracelet & \makecell[c]{-0.061\\{[-0.096, -0.028]}} & \makecell[c]{-0.126\\{[-0.179, -0.075]}} & \makecell[c]{-0.151\\{[-0.214, -0.088]}} & \makecell[c]{-0.116\\{[-0.192, -0.029]}} \\
\bottomrule
\end{tabular}
}
\caption*{\footnotesize Notes: Cells report adjusted predictions or Far-minus-Close
differences, with pointwise 95\% cluster-weight intervals in brackets. All
estimands retain all 9,805 randomized observations.}
\end{table}

\clearpage
\begin{table}[htbp]
\centering
\caption{Evolution of treatment-by-Far interaction effects}
\resizebox{\linewidth}{!}{\begin{tabular}{lcccc}
\toprule
 & Day 1 & Day 3 & Day 6 & Day 12 \\
\midrule
\multicolumn{5}{l}{\textit{Panel A: Arm-specific effects}} \\
Calendar $\times$ Far & \makecell[c]{0.003\\{[-0.039, 0.046]}} & \makecell[c]{-0.037\\{[-0.111, 0.035]}} & \makecell[c]{0.003\\{[-0.077, 0.083]}} & \makecell[c]{-0.010\\{[-0.110, 0.081]}} \\
Ink $\times$ Far & \makecell[c]{0.047\\{[0.004, 0.087]}} & \makecell[c]{0.071\\{[-0.002, 0.140]}} & \makecell[c]{0.143\\{[0.049, 0.237]}} & \makecell[c]{0.172\\{[0.048, 0.289]}} \\
Bracelet $\times$ Far & \makecell[c]{0.015\\{[-0.028, 0.059]}} & \makecell[c]{-0.010\\{[-0.080, 0.057]}} & \makecell[c]{0.024\\{[-0.056, 0.111]}} & \makecell[c]{0.065\\{[-0.042, 0.177]}} \\
\addlinespace
\multicolumn{5}{l}{\textit{Panel B: Pooled signal effect}} \\
Any Signal $\times$ Far & \makecell[c]{0.029\\{[0.001, 0.059]}} & \makecell[c]{0.049\\{[-0.002, 0.099]}} & \makecell[c]{0.082\\{[0.023, 0.143]}} & \makecell[c]{0.124\\{[0.052, 0.199]}} \\
\bottomrule
\end{tabular}
}
\caption*{\footnotesize Notes: Each arm-specific row is the treatment-versus-Control
effect under Far assignment minus the corresponding effect under Close assignment.
The pooled row compares the average of Ink and Bracelet with the average of Control
and Calendar. Brackets contain pointwise 95\% cluster-weight intervals.}
\end{table}

\begin{figure}[htbp]
\centering
\includegraphics[width=\linewidth]{figures/campaign-day-incidence.pdf}
\caption{Adjusted unconditional daily take-up. Shaded regions are pointwise
95\% cluster-weight intervals. These are unconditional daily probabilities, not
hazards among those who have not yet participated.}
\end{figure}

\paragraph{Summary.}
The pooled Any Signal--No Signal Far--Close interaction evolves as follows: day 1: 0.029 [0.001, 0.059]; day 3: 0.049 [-0.002, 0.099]; day 6: 0.082 [0.023, 0.143]; day 12: 0.124 [0.052, 0.199].
At day 12, the arm-specific treatment-by-Far interactions are Bracelet: 0.065 [-0.042, 0.177]; Calendar: -0.010 [-0.110, 0.081]; Ink: 0.172 [0.048, 0.289].
The day-12 estimates reproduce the existing original-assignment ATEs to a maximum absolute discrepancy of 1.7e-13.
Inference uses 1000 successful paired exponential cluster-weight draws out of 1000 attempted draws.


The differential effect of public signals under Far assignment therefore develops
over the campaign, rising from 2.9 percentage points on day 1 to 12.4 percentage
points by day 12. This accumulation is consistent with the value of visible signals
increasing as participation unfolds within communities, and provides additional
support for the broad social-image or social-information mechanism. The timing
evidence does not, however, separately identify social image from related mechanisms
such as reminders, descriptive norms, or social learning. The arm-specific estimates
also show that the aggregate timing pattern is driven most clearly by Ink; the
Bracelet interaction moves in the same direction later in the campaign but remains
imprecisely estimated.

\end{document}
